\documentclass[11pt,a4paper]{article}
\usepackage[top=2.5cm,bottom=2.5cm,left=2.5cm,right=2.5cm]{geometry}
\usepackage{fontspec}
\usepackage{enumitem}
\usepackage{fancyhdr}
\usepackage{hyperref}
\usepackage{titlesec}
\usepackage{array}
\usepackage{tabularx}
\usepackage{setspace}

\setmainfont{Plus Jakarta Sans}[BoldFont={Plus Jakarta Sans Bold},ItalicFont={Plus Jakarta Sans Italic}]
\onehalfspacing
\tolerance=1
\emergencystretch=\maxdimen
\hyphenpenalty=10000
\hbadness=10000
\setlength{\parindent}{1cm}
\setlength{\parskip}{0pt}
\setlength{\tabcolsep}{4pt}

\titleformat{\section}{\fontsize{13}{16}\bfseries}{}{0em}{}
\titlespacing{\section}{0pt}{18pt}{6pt}
\titleformat{\subsection}{\fontsize{12}{15}\bfseries}{}{0em}{}
\titlespacing{\subsection}{0pt}{12pt}{4pt}
\titleformat{\subsubsection}{\fontsize{11}{14}\bfseries}{}{0em}{}
\titlespacing{\subsubsection}{0pt}{8pt}{3pt}

\pagestyle{fancy}
\fancyhf{}
\rhead{\small KerjaCUS! -- Business Requirements Document}
\rfoot{\small Halaman \thepage}
\lfoot{\small Dokumen Rahasia -- KerjaCUS! Platform}
\renewcommand{\headrulewidth}{0.4pt}
\renewcommand{\footrulewidth}{0.3pt}
\hypersetup{colorlinks=false,pdfborder={0 0 0}}

\newcolumntype{Y}{>{\centering\arraybackslash}X}

\begin{document}
\thispagestyle{fancy}
\begin{center}
\vspace*{3cm}
{\fontsize{15}{18}\bfseries BUSINESS REQUIREMENTS DOCUMENT\\(BRD)}
\vspace{1cm}
{\large Proyek: \textbf{KopiNusantara E-Commerce}}\\[0.3cm]
{\normalsize Platform e-commerce untuk penjualan kopi specialty Indonesia\\dengan fitur subscription, katalog produk, dan pembayaran terintegrasi}
\vspace{1cm}
{\normalsize Versi 1.0 \textbullet\ Bahasa Indonesia}
\vspace{0.5cm}
{\small Nomor Dokumen: BRD-KOPI-001-v1}
\end{center}
\vspace{2cm}
\noindent\small
\begin{tabularx}{\textwidth}{|X|X|}
\hline
\centering\textbf{Pemilik Dokumen}\arraybackslash & \centering\textbf{KerjaCUS! Platform}\arraybackslash \\\hline
Sponsor Bisnis & Ahmad Fadillah (Project Owner) \\\hline
Target Pengguna & Pecinta kopi, roastery owner, distributor \\\hline
Tujuan Utama & Membangun platform e-commerce kopi specialty dengan fitur subscription dan pembayaran online \\\hline
Status & Draft -- Menunggu Persetujuan \\\hline
Tanggal & 24 Maret 2026 \\\hline
Dibuat oleh & KerjaCUS! AI Scoping Engine \\\hline
\end{tabularx}\normalsize

\newpage

\section{A. Document Overview}
\noindent\small
\begin{tabularx}{\textwidth}{|X|X|}
\hline
\centering\textbf{Nama Proyek}\arraybackslash & \centering\textbf{KopiNusantara E-Commerce}\arraybackslash \\\hline
Jenis Dokumen & Business Requirements Document (BRD) \\\hline
Versi & 1.0 \\\hline
Tanggal & 24 Maret 2026 \\\hline
Stakeholder Utama & Ahmad Fadillah (Owner), Tim KerjaCUS! \\\hline
Tujuan Dokumen & Menjelaskan kebutuhan bisnis, ruang lingkup, dan metrik keberhasilan untuk platform KopiNusantara \\\hline
\end{tabularx}\normalsize

\section{B. Executive Summary}
KopiNusantara adalah platform e-commerce yang dirancang khusus untuk memfasilitasi penjualan kopi specialty dari berbagai daerah di Indonesia. Platform ini menghubungkan roastery lokal dengan konsumen yang mencari pengalaman kopi premium.

Pasar kopi specialty Indonesia tumbuh 15-20\% per tahun, namun mayoritas roastery kecil masih mengandalkan penjualan offline atau marketplace umum yang tidak menyediakan fitur khusus seperti subscription bulanan, profil rasa kopi, dan rekomendasi berbasis preferensi.

Platform ini akan menyediakan: katalog kopi dengan filter berdasarkan origin, roast level, dan flavor notes; sistem subscription bulanan; pembayaran terintegrasi via Midtrans (QRIS, VA, e-wallet); dan dashboard penjualan untuk roastery.

\section{C. Problem Statement}
Saat ini, roastery kopi specialty di Indonesia menghadapi beberapa tantangan utama:
\begin{itemize}[noitemsep,leftmargin=1.5cm]
\item Penjualan masih mengandalkan WhatsApp atau marketplace umum yang tidak mendukung fitur khusus kopi.
\item Tidak ada platform yang menyediakan profil rasa (flavor profile) untuk membantu konsumen memilih.
\item Subscription bulanan masih dikelola manual via spreadsheet.
\item Konsumen kesulitan menemukan roastery lokal berkualitas di luar kota besar.
\item Tidak ada sistem review dan rating yang terkurasi untuk kopi specialty.
\end{itemize}

\section{D. Business Objective}
\noindent\small
\begin{tabularx}{\textwidth}{|p{0.8cm}|X|X|X|}
\hline
\centering\textbf{No.}\arraybackslash & \centering\textbf{Tujuan Bisnis}\arraybackslash & \centering\textbf{Target}\arraybackslash & \centering\textbf{Indikator}\arraybackslash \\\hline
1 & Meningkatkan jangkauan pasar roastery & 50+ roastery terdaftar dalam 6 bulan & Jumlah roastery aktif di platform \\\hline
2 & Meningkatkan penjualan kopi specialty online & Rp 500 juta GMV di Q1 & Total transaksi bulanan \\\hline
3 & Membangun basis subscriber & 1.000+ subscriber aktif & Jumlah subscription aktif \\\hline
4 & Meningkatkan repeat purchase rate & 40\% repeat order & Persentase pelanggan yang order lebih dari 1x \\\hline
\end{tabularx}\normalsize

\section{E. Scope}
\noindent\small
\begin{tabularx}{\textwidth}{|X|X|}
\hline
\centering\textbf{In Scope}\arraybackslash & \centering\textbf{Out of Scope}\arraybackslash \\\hline
Katalog produk kopi dengan filter origin, roast level, flavor notes & Roasting service atau private label \\\hline
Sistem subscription bulanan (pilih paket, frekuensi) & Logistik dan pengiriman sendiri (pakai ekspedisi pihak ketiga) \\\hline
Pembayaran online via Midtrans (VA, QRIS, e-wallet) & Multi-currency (hanya Rupiah) \\\hline
Dashboard penjualan untuk roastery & Sistem POS untuk toko fisik \\\hline
Review dan rating produk & Social commerce / live selling \\\hline
Rekomendasi kopi berbasis preferensi & AI-powered flavor matching (fase selanjutnya) \\\hline
\end{tabularx}\normalsize

\section{F. Stakeholders (Roles and Responsibilities)}
\noindent\small
\begin{tabularx}{\textwidth}{|p{3.2cm}|p{2.8cm}|X|}
\hline
\centering\textbf{Stakeholder}\arraybackslash & \centering\textbf{Peran}\arraybackslash & \centering\textbf{Tanggung Jawab}\arraybackslash \\\hline
Ahmad Fadillah & Project Owner & Menetapkan visi produk, menyetujui BRD/PRD, dan memvalidasi hasil akhir \\\hline
Tim Developer (KerjaCUS!) & Development Team & Mengembangkan fitur sesuai spesifikasi PRD dan timeline \\\hline
Roastery Partner & Content Provider & Menyediakan data produk, foto, dan deskripsi kopi \\\hline
End User (Konsumen) & Target User & Memberikan feedback dan validasi fitur melalui UAT \\\hline
\end{tabularx}\normalsize

\section{G. Target Users}
\noindent\small
\begin{tabularx}{\textwidth}{|p{2.8cm}|X|X|}
\hline
\centering\textbf{Segmen}\arraybackslash & \centering\textbf{Kebutuhan Utama}\arraybackslash & \centering\textbf{Contoh Manfaat}\arraybackslash \\\hline
Pecinta Kopi & Menemukan kopi specialty berkualitas dari berbagai daerah & Eksplorasi rasa baru setiap bulan via subscription \\\hline
Roastery Owner & Platform penjualan online yang mudah dikelola & Dashboard penjualan, manajemen stok, analytics \\\hline
Gift Buyer & Opsi hadiah kopi premium & Paket gift dengan personalisasi pesan \\\hline
Office Manager & Kebutuhan kopi kantor berkualitas & Subscription kantor dengan volume discount \\\hline
\end{tabularx}\normalsize

\section{H. Business Needs and Requirements}
\begin{itemize}[noitemsep,leftmargin=1.5cm]
\item Sistem harus mendukung pendaftaran roastery dengan verifikasi dokumen usaha.
\item Katalog produk harus menampilkan informasi lengkap: origin, altitude, processing method, roast date, flavor notes.
\item Subscription harus fleksibel: pilih paket (250g/500g/1kg), frekuensi (mingguan/bulanan), dan preferensi rasa.
\item Pembayaran harus terintegrasi dengan minimal 3 metode: bank transfer, e-wallet, QRIS.
\item Review hanya bisa diberikan oleh pembeli terverifikasi.
\item Dashboard roastery harus menampilkan: penjualan harian, stok, pesanan masuk, dan analytics.
\item Platform harus responsive (mobile-first) karena 70\%+ traffic diperkirakan dari mobile.
\end{itemize}

\section{I. Business Rules}
\noindent\small
\begin{tabularx}{\textwidth}{|p{1.2cm}|X|}
\hline
\centering\textbf{Kode}\arraybackslash & \centering\textbf{Aturan Bisnis}\arraybackslash \\\hline
BR-01 & Setiap roastery harus memiliki minimal 3 produk aktif untuk bisa ditampilkan di platform. \\\hline
BR-02 & Subscription yang tidak dibayar dalam 3 hari otomatis di-cancel. \\\hline
BR-03 & Review hanya bisa diberikan setelah produk diterima (status delivered). \\\hline
BR-04 & Roastery harus memproses pesanan dalam 2 hari kerja, atau pesanan otomatis di-cancel dengan refund. \\\hline
BR-05 & Harga produk sudah termasuk packaging, belum termasuk ongkos kirim. \\\hline
\end{tabularx}\normalsize

\section{J. Expected Benefits}
\noindent\small
\begin{tabularx}{\textwidth}{|p{3.2cm}|X|}
\hline
\centering\textbf{Kategori}\arraybackslash & \centering\textbf{Uraian}\arraybackslash \\\hline
Akses Pasar & Roastery kecil mendapat akses ke pasar nasional tanpa investasi besar di marketing. \\\hline
Customer Experience & Konsumen mendapat pengalaman belanja kopi yang terkurasi dan informatif. \\\hline
Operasional & Subscription otomatis mengurangi beban operasional roastery dalam mengelola repeat order. \\\hline
Revenue & Platform mendapat komisi per transaksi dan biaya subscription management. \\\hline
Data & Insight tentang preferensi kopi konsumen Indonesia untuk pengembangan fitur selanjutnya. \\\hline
\end{tabularx}\normalsize

\section{K. Risks and Assumptions}
\noindent\small
\begin{tabularx}{\textwidth}{|p{4.2cm}|X|}
\hline
\centering\textbf{Risiko}\arraybackslash & \centering\textbf{Mitigasi / Asumsi}\arraybackslash \\\hline
Kualitas kopi tidak konsisten & Sistem rating dan review membantu filtering; roastery dengan rating rendah mendapat warning. \\\hline
Pengiriman merusak freshness & Panduan packaging untuk roastery; opsi same-day delivery untuk area Jabodetabek. \\\hline
Adopsi roastery lambat & Onboarding gratis 3 bulan pertama; tim sales aktif approach roastery. \\\hline
Kompetisi marketplace umum & Diferensiasi via fitur khusus kopi (flavor profile, subscription, cupping notes). \\\hline
\end{tabularx}\normalsize

\section{L. Success Metrics}
\noindent\small
\begin{tabularx}{\textwidth}{|p{3.2cm}|p{2cm}|p{2.2cm}|X|}
\hline
\centering\textbf{Metrik}\arraybackslash & \centering\textbf{Baseline}\arraybackslash & \centering\textbf{Target 6 Bln}\arraybackslash & \centering\textbf{Metode Ukur}\arraybackslash \\\hline
Jumlah roastery & 0 & 50+ & Dashboard admin \\\hline
GMV bulanan & Rp 0 & Rp 100 juta & Analytics transaksi \\\hline
Subscriber aktif & 0 & 1.000+ & Database subscription \\\hline
Repeat purchase & 0\% & 40\%+ & Cohort analysis \\\hline
Rating rata-rata & N/A & 4.5/5 & Aggregasi review \\\hline
\end{tabularx}\normalsize

\section{M. Constraints}
\begin{itemize}[noitemsep,leftmargin=1.5cm]
\item Budget pengembangan maksimal Rp 75 juta (termasuk BRD dan PRD).
\item Timeline MVP: 60 hari dari approval PRD.
\item Tim development: 1 backend, 1 frontend, 1 UI/UX designer (via KerjaCUS! matching).
\item Integrasi pembayaran hanya via Midtrans (sesuai kesepakatan).
\item Bahasa platform: Bahasa Indonesia (English di fase selanjutnya).
\end{itemize}

\section{N. High Level Timeline}
\noindent\small
\begin{tabularx}{\textwidth}{|p{3.2cm}|p{1.8cm}|X|p{2.8cm}|}
\hline
\centering\textbf{Fase}\arraybackslash & \centering\textbf{Durasi}\arraybackslash & \centering\textbf{Output Utama}\arraybackslash & \centering\textbf{PIC}\arraybackslash \\\hline
AI Scoping \& BRD & 1 minggu & BRD dokumen (ini) & KerjaCUS! AI \\\hline
PRD \& Tech Design & 1 minggu & PRD, wireframe, API design & KerjaCUS! AI + Owner \\\hline
Talent Matching & 3--5 hari & Tim development terbentuk & KerjaCUS! Platform \\\hline
Sprint 1: Core & 3 minggu & Katalog, auth, checkout & Development Team \\\hline
Sprint 2: Features & 3 minggu & Subscription, dashboard, review & Development Team \\\hline
Sprint 3: Polish & 2 minggu & Testing, bug fix, deployment & Development Team \\\hline
UAT \& Launch & 1 minggu & Go-live & Semua stakeholder \\\hline
\end{tabularx}\normalsize

\vspace{1.5cm}
\begin{center}\small
--- Akhir Dokumen ---\\
Dibuat secara otomatis oleh KerjaCUS! AI Scoping Engine\\
Versi 1.0 -- 24 Maret 2026
\end{center}

\end{document}
